\chapter{Quantization and Deployment}
\label{ch:quant}

After training, the model is compressed so that it runs faster and uses less
memory. NVIDIA applies \emph{post-training quantization} (PTQ) to the weights
and activations of the post-trained BF16 model and releases two deployment
checkpoints: an FP8 checkpoint (8-bit weights and activations) for
Hopper-generation GPUs, and an NVFP4 checkpoint (4-bit weights and
activations) for the newer Blackwell GPUs~\cite{nvidia2026super}.

\section{FP8 Checkpoint}

The FP8 checkpoint was calibrated on 256 samples of 65{,}536 tokens each,
drawn from the SFT data. Quantization is applied to the MoE matrix
multiplications (both routed and shared experts) and to the Mamba linear
layers, where most of the computation is. The KV cache is kept in FP8, and the
Mamba state cache is stored in FP16 rather than FP32 for speed.
Table~\ref{tab:precision-fp8} lists the assignments.

\begin{table}[htbp]
\centering
\small
\caption[Precision settings of the FP8 checkpoint]{Precision settings of the
FP8 checkpoint and of the BF16 baseline (from~\cite{nvidia2026super},
Table~6)}
\label{tab:precision-fp8}
\begin{tabular}{@{}lcc@{}}
\toprule
\textbf{Component} & \textbf{FP8 checkpoint} & \textbf{BF16 baseline} \\
\midrule
Embedding & BF16 & BF16 \\
Attention projections (QKV and output) & BF16 & BF16 \\
KV cache and first attention product & FP8 & FP8 \\
Second attention product & BF16 & BF16 \\
\midrule
MoE experts (routed and shared) & FP8 & BF16 \\
MoE latent projections & BF16 & BF16 \\
Router & FP32 & FP32 \\
\midrule
Mamba linear layers & FP8 & BF16 \\
Mamba SSM cache & FP16 & FP32 \\
Mamba 1-D convolution & BF16 & BF16 \\
\midrule
Output layers & BF16 & BF16 \\
\bottomrule
\end{tabular}
\end{table}

\section{NVFP4 Checkpoint}

FP4 is a more aggressive format than FP8. It is particularly attractive for
prefill-heavy workloads such as coding agents, in which very long inputs are
processed and the linear and MoE matrix multiplications dominate the running
time. NVFP4 for inference uses the same element format as in training: signed
E2M1 values with a scale per block of 16, the block scales being stored in FP8
with a per-tensor FP32 factor.

Naive 4-bit quantization costs accuracy. The final recipe combines three
measures:
\begin{enumerate}
  \item \textbf{MSE-calibrated weight scales.} In the default recipe each
        block scale is set from the largest absolute value in the block.
        Because weights are quantized once, offline, an expensive search can
        be afforded instead, choosing the scale that minimises the mean squared
        error of the block.
  \item \textbf{Max-based activation scales.} Activations must be quantized at
        run time, where a search is impractical, so they keep the simple
        max-based scaling.
  \item \textbf{AutoQuantize.} Sensitive layers are selectively promoted from
        FP4 to FP8 or BF16.
\end{enumerate}

\emph{AutoQuantize}, part of NVIDIA's Model-Optimizer toolkit, selects the
best precision for each layer automatically. It estimates how sensitive each
operation is to quantization, using a second-order approximation of the change
in loss, models the performance cost of each candidate precision, and then
solves a knapsack-style optimisation for the assignment that loses the least
accuracy within a given budget. For \NS\ the candidates were NVFP4, FP8 and
BF16 with an effective budget of 4.75 bits. The outcome
(Table~\ref{tab:precision-fp4}) is a model that is mostly FP4: the routed
experts, which hold the great majority of the parameters, are NVFP4
throughout, while the attention and Mamba projections were assigned FP8 or
BF16.

\begin{table}[htbp]
\centering
\small
\caption[Precision settings of the NVFP4 checkpoint]{Precision settings of the
NVFP4 checkpoint (from~\cite{nvidia2026super}, Table~7)}
\label{tab:precision-fp4}
\begin{tabular}{@{}lcc@{}}
\toprule
\textbf{Component} & \textbf{Searched?} & \textbf{NVFP4 checkpoint} \\
\midrule
Embedding & No & BF16 \\
Attention QKV projection & Yes & BF16 \\
Attention output projection & Yes & FP8 / BF16 \\
KV cache and first attention product & No & FP8 \\
\midrule
Routed (sparse) experts & Yes & NVFP4 \\
Shared expert & Yes & NVFP4 / FP8 / BF16 \\
MoE latent projections & Yes & FP8 / BF16 \\
Router & No & FP32 \\
\midrule
Mamba projections & Yes & FP8 / BF16 \\
Mamba 1-D convolution & No & BF16 \\
Mamba SSM cache & No & FP16 \\
\midrule
Output layers & No & BF16 \\
\bottomrule
\end{tabular}
\end{table}

The whole mixed-precision PTQ process took less than two hours on a single
node with eight B200 GPUs, using 512 calibration samples of 4096 tokens. The
resulting model achieves a median of 99.8\% of the BF16 baseline's accuracy
while keeping near-FP4 speed. Table~\ref{tab:quant-eval} compares the three
checkpoints on a selection of benchmarks; the differences are small and go in
both directions.

\begin{table}[htbp]
\centering
\small
\caption[Accuracy of the BF16, FP8 and NVFP4 checkpoints]{Accuracy of the
BF16, FP8 and NVFP4 checkpoints on selected benchmarks
(from~\cite{nvidia2026super}, Table~8)}
\label{tab:quant-eval}
\begin{tabular}{@{}lccc@{}}
\toprule
\textbf{Benchmark} & \textbf{BF16} & \textbf{FP8} & \textbf{NVFP4} \\
\midrule
MMLU-Pro & 83.73 & 83.63 & 83.33 \\
HMMT Feb25 (with tools) & 94.73 & 94.38 & 95.36 \\
GPQA (no tools) & 79.23 & 79.36 & 79.42 \\
LiveCodeBench v5 & 81.19 & 80.99 & 80.56 \\
SciCode (subtask) & 42.05 & 41.38 & 40.83 \\
Terminal Bench (hard subset) & 25.78 & 26.04 & 24.48 \\
SWE-Bench (OpenCode) & 60.47 & --- & 59.90 \\
TauBench V2 (average) & 61.15 & 61.07 & 60.46 \\
IFBench (prompt) & 72.58 & 72.32 & 73.30 \\
AA-LCR & 58.31 & 57.69 & 58.06 \\
RULER 512k & 95.22 & 95.15 & 95.21 \\
RULER 1M & 91.64 & 91.43 & 91.60 \\
MMLU-ProX (average) & 79.35 & 79.21 & 79.37 \\
\bottomrule
\end{tabular}
\end{table}

\section{Mamba State Quantization}

\subsection{The Problem}

The most instructive part of the deployment work concerns the Mamba state. In
memory-bound settings, reading the Mamba state cache (the \emph{SSM cache})
from memory becomes a major bottleneck for decoding speed. The cache is stored
in FP32 by default. The obvious optimisation is to store it in FP16: fetch it
in FP16, convert to FP32 for the recurrent update, and convert back to FP16
for storage.

On an early checkpoint this simple change had an unexpected effect. Accuracy
on coding benchmarks hardly moved, but the model's answers became much
longer: the number of generated tokens on LiveCodeBench rose by about 37\%
with BF16 weights and by about 40\% with FP8 weights
(Table~\ref{tab:ssm-cache}). A model that writes 40\% more tokens to reach the
same answer has given back much of the speed that quantization was meant to
provide.

\begin{table}[htbp]
\centering
\small
\caption[Effect of the SSM cache recipe on accuracy and verbosity]{Effect of
the SSM cache recipe on LiveCodeBench accuracy and verbosity (selected rows
from~\cite{nvidia2026super}, Table~9). W16A16 denotes BF16 weights and
activations, W8A8 denotes FP8.}
\label{tab:ssm-cache}
\begin{tabular}{@{}llccc@{}}
\toprule
\textbf{Weights/act.} & \textbf{SSM cache recipe} & \textbf{Accuracy} &
\textbf{Completion tokens} & \textbf{Verbosity change} \\
\midrule
W16A16 & FP32 (baseline) & 72.91 & 21{,}769 & 0.00\% \\
W16A16 & FP16 & 73.24 & 29{,}812 & $+36.95\%$ \\
W16A16 & FP16 + SR (Philox 5) & 72.00 & 21{,}392 & $-1.73\%$ \\
\midrule
W8A8 & FP16 & 73.13 & 30{,}536 & $+40.27\%$ \\
W8A8 & INT16 + block scale (128) & 72.22 & 22{,}406 & $+2.90\%$ \\
W8A8 & FP16 + SR (Philox 5) & 72.63 & 22{,}159 & $+1.79\%$ \\
W8A8 & FP16 + SR (Philox 3) & 70.07 & 24{,}098 & $+10.70\%$ \\
\bottomrule
\end{tabular}
\end{table}

\subsection{Why the Error Accumulates}

The cause lies in the recurrence itself. In an attention layer, a rounding
error in one cache entry stays in that entry. In a Mamba layer the state is
fed back into itself at every step, so an error made at one step is carried
into all later steps. Let the exact update be $h_t = A_t h_{t-1} + B_t x_t$ and
let storing the state at step $t$ add an error $e_t$, so that the quantized
state follows $h_{q,t} = A_t h_{q,t-1} + B_t x_t + e_t$. Unrolling the
recursion gives~\cite{nvidia2026super}
\begin{equation}
  h_{q,t} = h_t + \sum_{i=0}^{t} \left( \prod_{j=i+1}^{t} A_j \right) e_i .
  \label{eq:error-accum}
\end{equation}
Every past error $e_i$ is still present at step $t$, scaled by the product of
the transitions since. Over the tens of thousands of steps in a long
generation, small errors can add up to a visible change in behaviour.

Fixing this through training is not straightforward. Mamba-2 is trained with
the chunked SSD algorithm (Section~\ref{sec:mamba2}), which never
materialises the step-by-step recurrence or the inference-time cache, so
simulating cache quantization during training would add substantial overhead.
The team therefore looked for training-free fixes.

\subsection{Two Fixes}

\textbf{More mantissa bits.} Storing the cache as 16-bit integers gives more
precision than FP16 for values of similar size. Plain INT16 did not help,
because the state has a wide dynamic range, but adding an FP32 scale per block
of 128 values along the state dimension removed the verbosity problem.

\textbf{Stochastic rounding.} The second hypothesis was that the trouble comes
from \emph{how} values are rounded. The standard rule, round-to-nearest with
ties to even (RTNE), always maps the same input to the same output. Its error
is therefore a systematic \emph{bias}, and in a recurrence a bias accumulates
coherently, pushing the state steadily in one direction. Stochastic rounding
(SR) instead rounds a value $x$ lying between two representable neighbours
$x_{\downarrow} \le x \le x_{\uparrow}$ at random, with probabilities
proportional to how close each neighbour is:
\begin{equation}
  \operatorname{SR}(x) =
  \begin{cases}
    x_{\uparrow}   & \text{with probability } p = \dfrac{x - x_{\downarrow}}{x_{\uparrow} - x_{\downarrow}},\\[8pt]
    x_{\downarrow} & \text{with probability } 1 - p ,
  \end{cases}
  \qquad \mathbb{E}\bigl[\operatorname{SR}(x)\bigr] = x .
  \label{eq:sr}
\end{equation}
Each individual rounding is noisier, but the error has zero mean, so the
errors in Eq.~\eqref{eq:error-accum} tend to cancel instead of adding up: the
systematic drift is replaced by zero-mean noise. Applying stochastic rounding
before the cast to FP16 fixed the verbosity problem for both the BF16 and the
FP8 models.

\subsection{The Chosen Recipe}

Both fixes restored accuracy and verbosity to the FP32 baseline. \NS\ uses
FP16 with stochastic rounding, for three reasons given in the report: it needs
no block scale factors to be computed, stored and loaded; Blackwell GPUs
provide a dedicated instruction for stochastic rounding during type
conversion; and they support the Philox pseudo-random generator natively. The
number of Philox rounds trades randomness quality against overhead. Five
rounds kept accuracy and verbosity at baseline, while three rounds were too
few (verbosity $+10.7\%$), so Philox with five rounds was selected.

This episode is a good illustration of a general point about hybrid models: a
state-space layer is not simply a cheaper attention layer. Its recurrent
nature creates failure modes that Transformers do not have, and the
deployment stack has to be adapted accordingly.
